
\chapterrule
\section{怎样让损失变小？梯度下降调参数}
\label{intro:16a-gradient-descent}

上一章比较了几组参数，却还没有找到最佳直线。
这个三点直线拟合问题可以直接求出最优参数；如果只要这两个数，直接解更简便。
后面的多层网络却通常很难一次解出全部最优参数，需要一种能根据当前误差反复调整参数的方法。
本章先求出这个小例子的最优解，再把它作为已知终点，核验\textbf{逐步更新}能否走近它。

现在从候选 A 的 $(w,b)=(1,0)$ 出发，先只观察 $w$ 变化时损失怎样变化，
再求出最佳参数、练习同步更新。全程使用同一组三点数据。

\subsection{先试一点：哪一边让损失下降？}

保持 $b=0$，暂时只改变 $w$。
把 $w=0.9,1,1.1$ 代入上一章按残差写出的损失，得到
\[
L(0.9,0)=1.74,\qquad L(1,0)=\frac53,\qquad L(1.1,0)\approx1.6067.
\]
在当前位置附近，增大 $w$ 能降低损失。
如果只靠逐个尝试，每个参数各选 $10$ 个值，两个参数的所有组合就要比较 $10\times10=100$ 次；
参数更多时，组合数还会继续相乘。
真实模型可有上亿个参数，
穷举这样的组合难以承担，而且有限的候选还可能漏掉更好的值。
我们可以用\textbf{导数}描述这样的局部变化率，而不必每次试两个方向。
相关符号见\hyperref[app:math-tools]{数学速查}。

固定 $b=0$ 时，同一个损失为
\[
L(w,0)=\frac{(1-w)^2+4+(w-2)^2}{3}.
\]
分别对三个平方项求导：
\[
\frac{\mathrm dL(w,0)}{\mathrm dw}
=\frac{-2(1-w)+2(w-2)}{3}
=\frac43w-2.
\]
在 $w=1$ 处，导数为 $-2/3$：增大一点 $w$，损失会减小，与刚才的试算一致。
导数的严格定义为
\[
\frac{\mathrm dL}{\mathrm dw}
=\lim_{\varepsilon\to0}\frac{L(w+\varepsilon)-L(w)}{\varepsilon}.
\]
这里保持 $b$ 不变，比较 $w$ 的微小变化引起的损失变化。
图~\ref{fig:loss-descent} 把三个试算点与当前点的切线画在一起：
附近的曲线向右下方延伸，切线的负斜率概括了这一局部趋势。
% Fixed-b=1 slice of the same three-point MSE used for the joint w/b update.
\par\noindent\begin{minipage}{\linewidth}
\tagstructbegin{tag=Figure,alttext={固定偏置 b 等于一时，一维损失曲线从 w 等于一的损失三分之二，下降到 w 等于一点四的零点二四，再到一点六四的零点零八六四，最优斜率为二}}
\begin{center}
\begin{tikzpicture}
\begin{axis}[
 width=118mm,height=56mm,
 xmin=.78,xmax=2.30,ymin=-.05,ymax=.86,
 axis lines=left,xlabel={$w$},ylabel={$L(w,1)$},
 xtick={1,1.4,1.64,2},ytick={0,.24,.6667},
 tick label style={font=\footnotesize},
 label style={font=\small},clip=false,
]
\addplot[datacolor,very thick,domain=.8:2.3,samples=80] {2/3*(x-2)^2};
\addplot[datacolor,only marks,mark=*,mark size=2.8pt]
 coordinates {(1,.6666667) (1.4,.24) (1.64,.0864)};
\addplot[datacolor,only marks,mark=*,mark size=2.5pt] coordinates {(2,0)};
\draw[-{Latex},datacolor,thick] (axis cs:1.05,.59) -- (axis cs:1.34,.29);
\draw[-{Latex},datacolor,thick] (axis cs:1.43,.225) -- (axis cs:1.61,.107);
\node[anchor=south,font=\footnotesize] at (axis cs:1,.68) {$w_0$};
\node[anchor=south,font=\footnotesize] at (axis cs:1.4,.25) {$w_1$};
\node[anchor=south west,font=\footnotesize] at (axis cs:1.66,.105) {$w_2$};
\node[anchor=north,font=\footnotesize,datacolor] at (axis cs:2,.0) {最优};
\end{axis}
\end{tikzpicture}

{\footnotesize 同一三点数据固定 $b=1$：$L(w,1)=\frac23(w-2)^2$；
三次参数位置的损失为 $2/3\to0.24\to0.0864$。}
\par\captionof{figure}{固定偏置时的一维梯度下降截面}\label{fig:loss-descent}\end{center}
\tagstructend
\end{minipage}\par


\subsection{两个参数各往哪里调？偏导数与梯度}
\phantomsection\label{concept:gradient}

实际拟合要调整\textbf{两个参数}。
\textbf{偏导数}在求某一个参数的变化率时，把其他参数视为常数。
对\eqref{eq:three-point-residual-loss} 的三个残差分别求导并相加，得到
\[
\begin{aligned}
\frac{\partial L}{\partial w}
&=\frac{2(-w+b+1)(-1)+2(w+b-2)}{3}=\frac43w-2,\\
\frac{\partial L}{\partial b}
&=\frac{2(-w+b+1)+2(b-2)+2(w+b-2)}{3}=2b-2.
\end{aligned}
\]
把两个偏导数排在一起，叫作\textbf{梯度}：
\[
\nabla L(w,b)=\left(\frac{\partial L}{\partial w},\frac{\partial L}{\partial b}\right)
=\left(\frac43w-2,2b-2\right).
\]
在初值 $(1,0)$ 处，$\nabla L(1,0)=\boxed{(-2/3,-2)}$。
两个分量均为负，分别提示增大一点 $w$、增大一点 $b$ 可以降低损失。
图~\ref{fig:gradient-partial-slices} 从同一个旧点分别观察这两个方向。
每幅图只改变一个参数；把两条切线的斜率排在一起，就是当前位置的梯度。
% Equal coordinate units in both panels; both pass through the old point (1,0).
\par\noindent\begin{minipage}{\linewidth}
\tagstructbegin{tag=Figure,alttext={在同一个参数点 w 等于一、b 等于零处观察两个局部截面。左图固定 b，只改变 w，切线斜率为负三分之二；右图固定 w，只改变 b，切线斜率为负二。两图的横轴范围均跨零点三、纵轴范围相同，右图切线更陡。两个斜率按 w、b 的顺序组成梯度。}}
\begin{center}
\begin{tikzpicture}
\begin{axis}[
 width=64mm,height=49mm,xmin=.85,xmax=1.15,ymin=1.30,ymax=2.04,
 axis lines=left,xlabel={$w$},ylabel={$L(w,0)$},
 title={固定 $b=0$},title style={font=\small},
 xtick={.9,1,1.1},ytick={1.4,1.666667,2},yticklabels={$1.4$,$5/3$,$2$},
 tick label style={font=\footnotesize},label style={font=\small},
]
\addplot[flowcolor,very thick,domain=.85:1.15,samples=50] {((1-x)^2+4+(x-2)^2)/3};
\addplot[residualcolor,thick,dashed,domain=.85:1.15] {5/3-2/3*(x-1)};
\addplot[datacolor,only marks,mark=*,mark size=2.4pt] coordinates {(1,1.666667)};
\node[font=\footnotesize] at (axis cs:1,1.93) {$\partial L/\partial w=-2/3$};
\end{axis}
\begin{axis}[
 at={(76mm,0)},width=64mm,height=49mm,xmin=-.15,xmax=.15,ymin=1.30,ymax=2.04,
 axis lines=left,xlabel={$b$},ylabel={$L(1,b)$},
 title={固定 $w=1$},title style={font=\small},
 xtick={-.1,0,.1},ytick={1.4,1.666667,2},yticklabels={$1.4$,$5/3$,$2$},
 tick label style={font=\footnotesize},label style={font=\small},
]
\addplot[flowcolor,very thick,domain=-.15:.15,samples=50] {x^2/3+(x-2)^2/3+(x-1)^2/3};
\addplot[residualcolor,thick,dashed,domain=-.15:.15] {5/3-2*x};
\addplot[datacolor,only marks,mark=*,mark size=2.4pt] coordinates {(0,1.666667)};
\node[font=\footnotesize] at (axis cs:.045,1.93) {$\partial L/\partial b=-2$};
\end{axis}
\end{tikzpicture}

{\footnotesize 实线是损失曲线，虚线是当前点的切线；两图使用相同的坐标单位。}
\captionof{figure}{偏导数每次只改变一个参数，梯度把同一旧点处的两个斜率合在一起}
\label{fig:gradient-partial-slices}
\end{center}
\tagstructend
\end{minipage}\par


\subsection{先直接求解：最优参数在哪里？}
\label{intro:gradient-direct-solution}

刚才已经算出两个偏导数。令它们同时为零，只需解两条一次方程：
\[
\begin{aligned}
\frac43w-2=0&\quad\Longrightarrow\quad w=\frac32,\\
2b-2=0&\quad\Longrightarrow\quad b=1.
\end{aligned}
\]
代回模型，三项预测为 $(-0.5,1,2.5)$，残差为 $(0.5,-1,0.5)$，因此
\[
L\left(\frac32,1\right)=\frac{0.5^2+(-1)^2+0.5^2}{3}=\boxed{\frac12}.
\]
本例的\textbf{最优参数为 $(3/2,1)$，最小损失为 $1/2$}。
为什么这确实是全部参数中的最小值，见后面的\hyperref[proof:three-point-global-minimum]{选读证明框}。

\begin{boundarybox}[title={能直接解，为什么还要学迭代？}]
本例的两个方程彼此独立，直接求解很容易。
一般多层网络中，许多参数共同影响输出，令梯度为零会得到相互耦合且含非线性的方程组，
通常很难直接求解；即使找到了梯度为零的点，也未必是全局最小值。
梯度下降提供另一条路：每一步计算当前位置的梯度，再按同一规则调整参数。
这个小例子的终点已知，正适合检查这条路怎样走。
\end{boundarybox}

\subsection{沿负梯度走一步，再检查新预测}
\phantomsection\label{concept:gradient-descent}

终点 $(3/2,1)$ 我们已经知道了。现在从 $(1,0)$ 出发，
更新时只用当前位置的梯度，看能否逐步走近它。
\textbf{梯度下降}沿梯度的反方向更新参数，以\textbf{学习率} $\eta$ 控制步长。
两个参数都用同一组旧参数处的梯度计算，这叫\textbf{同步更新}：
\begin{equation}
\boxed{\begin{aligned}
w_{\rm new}&=w_{\rm old}-\eta\frac{\partial L}{\partial w},\\
b_{\rm new}&=b_{\rm old}-\eta\frac{\partial L}{\partial b}.
\end{aligned}}
\label{eq:gradient-simple}
\end{equation}
取 $\eta=0.3$，从 $(1,0)$ 出发，第一步为
\[
(w_1,b_1)=(1,0)-0.3(-2/3,-2)=\boxed{(1.2,0.6)}.
\]
新预测为 $(-0.6,0.6,1.8)$，与真实值 $(-1,2,2)$ 比较，残差为 $(0.4,-1.4,-0.2)$，所以
\[
L(1.2,0.6)=\frac{0.4^2+(-1.4)^2+(-0.2)^2}{3}=\boxed{0.72}<\frac53.
\]
这一步改善了初始预测。它还没胜过上一章候选 B 的 $2/3$，可以继续按同一规则更新。
图~\ref{fig:gradient-prediction-update} 回到样本的 $x,y$ 平面，比较更新前后的直线与残差。
左端样本原先恰好猜对，更新后出现了 $0.4$ 的残差；
其余两项的改善使整体损失下降。\textbf{整体变好，不要求每个样本都变好。}
% Same data-space limits and marker meanings as the optimal-fit panel.
\par\noindent\begin{minipage}{\linewidth}
\tagstructbegin{tag=Figure,alttext={用同一个旧点处的梯度负三分之二、负二和学习率零点三，将参数从一、零同步更新为一点二、零点六。左右两图分别画更新前后的直线，黑实点为真实样本，空心方点为预测，竖线为残差。更新前残差为零、负二、负一，损失三分之五；更新后残差为零点四、负一点四、负零点二，损失零点七二。左端样本变差，但整组预测的损失下降。}}
\begin{center}
{\small $(1,0)\xrightarrow{\eta=0.3,\ \nabla L=(-2/3,-2)}(1.2,0.6)$}\par\smallskip
\begin{tikzpicture}
\begin{axis}[
 name=before,width=64mm,height=54mm,xmin=-1.25,xmax=1.25,ymin=-1.5,ymax=3.1,
 axis lines=left,xlabel={$x$},ylabel={$y$},title={更新前：$\widehat{y}=x$},
 title style={font=\small},xtick={-1,0,1},ytick={-1,0,1,2,3},
 tick label style={font=\footnotesize},label style={font=\small},
]
\addplot[flowcolor,thick,dashed,domain=-1.15:1.15] {x};
\draw[residualcolor,thick] (axis cs:-1,-1)--(axis cs:-1,-1);
\draw[residualcolor,thick] (axis cs:0,2)--(axis cs:0,0);
\draw[residualcolor,thick] (axis cs:1,2)--(axis cs:1,1);
\addplot[only marks,datacolor,mark=square*,mark options={fill=white},mark size=2.4pt]
 coordinates {(-1,-1) (0,0) (1,1)};
\addplot[only marks,noisecolor,mark=*,mark size=1.7pt] coordinates {(-1,-1) (0,2) (1,2)};
\end{axis}
\begin{axis}[
 name=after,at={(76mm,0)},width=64mm,height=54mm,xmin=-1.25,xmax=1.25,ymin=-1.5,ymax=3.1,
 axis lines=left,xlabel={$x$},ylabel={$y$},title={更新后：$\widehat{y}=1.2x+0.6$},
 title style={font=\small},xtick={-1,0,1},ytick={-1,0,1,2,3},
 tick label style={font=\footnotesize},label style={font=\small},
]
\addplot[flowcolor,thick,domain=-1.15:1.15] {1.2*x+.6};
\draw[residualcolor,thick] (axis cs:-1,-1)--(axis cs:-1,-.6);
\draw[residualcolor,thick] (axis cs:0,2)--(axis cs:0,.6);
\draw[residualcolor,thick] (axis cs:1,2)--(axis cs:1,1.8);
\addplot[only marks,datacolor,mark=square*,mark options={fill=white},mark size=2.4pt]
 coordinates {(-1,-.6) (0,.6) (1,1.8)};
\addplot[only marks,noisecolor,mark=*,mark size=1.7pt] coordinates {(-1,-1) (0,2) (1,2)};
\end{axis}
\node[anchor=north,font=\footnotesize,align=center] at ([yshift=-2mm]before.outer south)
 {残差 $(0,-2,-1)$\\$L=5/3$};
\node[anchor=north,font=\footnotesize,align=center] at ([yshift=-2mm]after.outer south)
 {残差 $(0.4,-1.4,-0.2)$\\$L=0.72$};
\end{tikzpicture}

{\footnotesize 黑实点：真实值；空心方点：预测；竖线：残差。两图坐标范围相同。}
\captionof{figure}{一次同步更新改变了整条预测直线，整体损失下降而各点改善程度不同}
\label{fig:gradient-prediction-update}
\end{center}
\tagstructend
\end{minipage}\par


\subsection{连走五步：每次都重新求梯度}
\phantomsection\label{experiment:beginner-iteration}

继续使用 $\eta=0.3$。每行先算该行参数的梯度，再得到下一行；
两个参数都从同一行的旧值出发。
\begin{center}\small
\tagpdfsetup{table/header-rows={1}}
\begin{tabular}{rrrr}
\toprule
已更新步数 $t$ & $w_t$ & $b_t$ & $L(w_t,b_t)$\\\midrule
0 & 1.000000 & 0.000000 & 1.666667\\
1 & 1.200000 & 0.600000 & 0.720000\\
2 & 1.320000 & 0.840000 & 0.547200\\
3 & 1.392000 & 0.936000 & 0.511872\\
4 & 1.435200 & 0.974400 & 0.503455\\
5 & 1.461120 & 0.989760 & 0.501113\\
10 & 1.496977 & 0.999895 & 0.500006\\
\bottomrule
\end{tabular}
\end{center}
例如更新一步后的梯度为 $(-0.4,-0.8)$，因此第二步得到
\[
(1.2,0.6)-0.3(-0.4,-0.8)=(1.32,0.84).
\]
第二步的损失 $0.5472$ 已低于候选 B 的 $2/3$。
继续更新时，参数逐渐接近已知的 $(1.5,1)$，损失也逐渐接近已知的最小值 $0.5$。
这些数值验证了本例在所选学习率下的迭代效果。
运行下列命令可核对整张表：
\begin{verbatim}
python3 scripts/beginner_training_demo.py
\end{verbatim}
图~\ref{fig:gradient-iteration-path} 用两种坐标记录同一段迭代。
左图的每个点代表一组参数 $(w,b)$，区别于上一图中代表样本的 $(x,y)$；
右图把这组参数对应的损失按更新步数画出。每一步重新求梯度，得到的移动方向和距离也随之变化。
% Joint full-batch updates, not the earlier fixed-b local observation.
\par\noindent\begin{minipage}{\linewidth}
\tagstructbegin{tag=Figure,alttext={同一段学习率零点三的迭代，左图画在参数平面，横轴为 w、纵轴为 b，每点是一组参数。轨迹从一、零经过一点二、零点六及一点三二、零点八四，逐渐接近前面已求出的最优参数一点五、一。右图横轴为更新步数，纵轴为整组三点的损失，标出相同的第零、一、二、五、十步；第二步损失零点五四七二已低于候选 B 的三分之二，随后趋近点线所示的已知最小损失二分之一。}}
\begin{center}
\begin{tikzpicture}
\begin{axis}[
 width=64mm,height=57mm,axis equal image,
 xmin=.85,xmax=1.98,ymin=-.12,ymax=1.28,
 axis lines=left,xlabel={$w$},ylabel={$b$},title={参数轨迹},title style={font=\small},
 xtick={1,1.5},ytick={0,.5,1},
 tick label style={font=\footnotesize},label style={font=\small},clip=false,
]
\addplot[flowcolor,thick,mark=*,mark size=1.3pt] coordinates {
 (1,0) (1.2,.6) (1.32,.84) (1.392,.936) (1.4352,.9744)
 (1.46112,.98976) (1.476672,.995904) (1.4860032,.9983616)
 (1.49160192,.99934464) (1.494961152,.999737856) (1.4969766912,.9998951424)};
\addplot[datacolor,only marks,mark=*,mark size=2pt] coordinates {
 (1,0) (1.2,.6) (1.32,.84) (1.46112,.98976) (1.4969766912,.9998951424)};
\draw[-{Latex},flowcolor,thick,shorten <=2pt,shorten >=2pt] (axis cs:1,0)--(axis cs:1.2,.6);
\draw[-{Latex},flowcolor,thick,shorten <=2pt,shorten >=2pt] (axis cs:1.2,.6)--(axis cs:1.32,.84);
\node[anchor=west,font=\footnotesize] at (axis cs:1.04,.08) {$t=0$};
\node[anchor=west,font=\footnotesize] at (axis cs:1.26,.54) {$t=1$};
\node[anchor=east,font=\footnotesize] at (axis cs:1.27,.80) {$t=2$};
\draw[flowcolor,thin] (axis cs:1.46112,.98976)--(axis cs:1.65,.91)
 node[anchor=west,font=\footnotesize] {$t=5$};
\draw[flowcolor,thin] (axis cs:1.4969766912,.9998951424)--(axis cs:1.65,1.15)
 node[anchor=west,font=\footnotesize] {$t=10$};
\end{axis}
\begin{axis}[
 at={(76mm,0)},width=66mm,height=57mm,xmin=-.3,xmax=10.5,ymin=.42,ymax=1.82,
 axis lines=left,xlabel={更新步数 $t$},ylabel={$L(w_t,b_t)$},title={对应的损失},title style={font=\small},
 xtick={0,1,2,5,10},ytick={.5,.6666667,1,1.6666667},yticklabels={$0.5$,$2/3$,$1$,$5/3$},
 tick label style={font=\footnotesize},label style={font=\small},
]
\addplot[noisecolor,dashed,domain=0:10] {2/3};
\addplot[noisecolor,densely dotted,thick,domain=0:10] {1/2};
\addplot[flowcolor,thick,mark=*,mark size=1.3pt] coordinates {
 (0,1.6666666667) (1,.72) (2,.5472) (3,.511872) (4,.50345472)
 (5,.5011126272) (6,.5003795743) (7,.5001332913) (8,.5000474480)
 (9,.5000169954) (10,.5000061046)};
\addplot[datacolor,only marks,mark=*,mark size=2pt] coordinates {
 (0,1.6666666667) (1,.72) (2,.5472) (5,.5011126272) (10,.5000061046)};
\node[font=\footnotesize] at (axis cs:6.8,.83) {候选 B：$L=2/3$};
\end{axis}
\end{tikzpicture}

{\footnotesize 右图虚线为候选 B 的 $L=2/3$，点线为已知最小损失 $L=1/2$。}
\captionof{figure}{本例取 $\eta=0.3$ 连续更新：参数逐步接近已知最优点，损失趋近已知最小值}
\label{fig:gradient-iteration-path}
\end{center}
\tagstructend
\end{minipage}\par


图~\ref{fig:loss-two-parameter-contours} 再把已知最优点与迭代路径放在同一张图中：
左图中，同一条等高线上的参数具有相同损失，迭代从外侧走向损失更低的内侧；
右图中，最优直线仍留下三个非零残差，与最小损失 $1/2$ 一致。
% Contours are exact for L(w,b)=2/3(w-2)^2+(b-1)^2.
% Radii correspond to contour levels 5/3 and 0.4.
\tagstructbegin{tag=Figure,alttext={双参数均方误差的椭圆等高线，中心为最优解 w 等于二 b 等于一，初始参数一和零沿梯度方向更新至一点四和零点六，损失由三分之五降低到零点四}}
\begin{center}
\begin{tikzpicture}[x=22mm,y=16mm,>=Latex]
\draw[->,gray] (0,-.6)--(4,-.6) node[right] {$w$};
\draw[->,gray] (0,-.6)--(0,2.6) node[above] {$b$};
\draw[blue,thick] (2,1) ellipse [x radius=34.79mm,y radius=20.66mm];
\draw[teal,thick] (2,1) ellipse [x radius=17.04mm,y radius=10.12mm];
\node[blue,font=\scriptsize] at (3.58,1.6) {$L=5/3$};
\node[teal,font=\scriptsize] at (2.84,1.15) {$L=0.4$};
\fill[orange] (1,0) circle (2.5pt) node[below,font=\scriptsize] {$(1,0)$};
\fill[orange] (1.4,.6) circle (2.5pt) node[below right,font=\scriptsize] {$(1.4,0.6)$};
\fill[teal] (2,1) circle (2.5pt) node[above,font=\scriptsize] {$(2,1)$ 最优};
\draw[->,orange,very thick] (1.04,.06)--(1.36,.54);
\end{tikzpicture}
\end{center}

\begin{center}
\captionof{figure}{同一三点 MSE 的双参数等高线：更新 $w$ 与 $b$ 后向低损失区域移动}
\label{fig:loss-two-parameter-contours}
\end{center}
\tagstructend


\subsection{一次训练要做哪四件事？}
\label{intro:training-four-steps}

把刚才每次更新的过程归纳起来，就是反复执行四件事：
\begin{enumerate}
\item \textbf{前向传播}：用当前参数计算 $\widehat{y}$。
\item \textbf{计算损失}：比较预测 $\widehat{y}$ 与真实目标 $y$。
\item \textbf{求梯度}：求损失对各个参数的偏导数。
\item \textbf{更新参数}：统一执行 $\theta\leftarrow\theta-\eta\nabla_\theta L$。
\end{enumerate}
这里 $\theta$ 统称需要训练的参数，本例就是 $w,b$。
更新后回到第一件事，用新参数重新预测、计算损失和梯度。
后面的网络会增加前向传播与求梯度的内部步骤，这四件事仍构成训练的基本循环。\phantomsection\label{route:gradient-end}

\subsection{学习率多大才稳？把“会发散”算出来}
\phantomsection\label{experiment:learning-rate-stability}

方向有了，步子太大会怎样？仍从 $(1,0)$ 出发，把学习率改为 $\eta=1.2$。
只看 $b$，更新式是 $b_{t+1}=b_t-1.2(2b_t-2)$，例如第一步为 $0-1.2\times(-2)=2.4$。
为了看清它是否靠近前面求出的最优值 $b=1$，继续更新，并记录每一步的差值 $b_t-1$：
\begin{center}\small
\tagpdfsetup{table/header-rows={1}}
\begin{tabular}{rrrr}
\toprule
已更新步数 $t$ & $b_t$ & 偏差 $b_t-1$ & 距离 $|b_t-1|$\\\midrule
0 & 0 & $-1$ & 1\\
1 & 2.4 & 1.4 & 1.4\\
2 & $-0.96$ & $-1.96$ & 1.96\\
3 & 3.744 & 2.744 & 2.744\\
\bottomrule
\end{tabular}
\end{center}
这里把带正负号的差值 $b_t-1$ 叫作\textbf{参数偏差}。
例如 $b_0=0$ 时，偏差为 $0-1=-1$，表示在最优值左边，距离为 $1$；
$b_1=2.4$ 时，偏差为 $2.4-1=1.4$，表示在最优值右边，距离为 $1.4$。
所以，\textbf{偏差的正负表示在哪一边，绝对值表示距离。}
表中 $b$ 在最优值两侧来回跳，距离却逐步变大；这就是这个方向上的\textbf{发散}。

为什么这些距离恰好每步变成原来的 $1.4$ 倍？
我们要找出\textbf{更新后的偏差与更新前的偏差有什么关系}。
先把学习率写回一般的 $\eta$，并把梯度 $2b_t-2$ 写成 $2(b_t-1)$：
\[
b_{t+1}=b_t-\eta(2b_t-2)=b_t-2\eta(b_t-1).
\]
左边的 $b_{t+1}$ 是更新后的参数。要得到更新后的偏差 $b_{t+1}-1$，
就在\textbf{等式两边同时减去最优值 $1$}：
\[
\begin{aligned}
b_{t+1}-1
&=b_t-2\eta(b_t-1)-1\\
&=(b_t-1)-2\eta(b_t-1)\\
&=(1-2\eta)(b_t-1).
\end{aligned}
\]
第二行把 $b_t$ 与末尾的 $-1$ 合在一起；第三行提取两项共有的 $b_t-1$。
这个结果说的是：\textbf{每更新一步，偏差就乘一次 $1-2\eta$。}
这里同时减去 $1$ 是为了分析偏差怎样变化；训练仍按原来的更新式执行。

当 $\eta=1.2$ 时，这个乘数为 $1-2\times1.2=-1.4$，因此偏差依次变为
\[
-1\ \xrightarrow{\times(-1.4)}\ 1.4
\ \xrightarrow{\times(-1.4)}\ -1.96
\ \xrightarrow{\times(-1.4)}\ 2.744.
\]
负号让偏差每步换符号，也就是跳到最优值的另一侧；
乘数的绝对值 $1.4$ 则让距离每步变成原来的 $1.4$ 倍，即增加 $40\%$。
把这个控制偏差变化的乘数记为 $r_b=1-2\eta$。

对 $w$ 也可以这样分析：它的最优值是 $3/2$，所以偏差为 $w_t-3/2$。
把梯度 $\frac43w_t-2$ 写成 $\frac43(w_t-\frac32)$，再在更新式两边同时减去 $3/2$，得到
\[
w_{t+1}-\frac32=\left(1-\frac43\eta\right)\left(w_t-\frac32\right),
\qquad r_w=1-\frac43\eta.
\]
取 $\eta=0.3$ 时，$r_w=0.6,r_b=0.4$：
初始偏差 $(-0.5,-1)$ 变成 $(-0.3,-0.4)$，正好得到 $(1.2,0.6)$。
第 $t$ 步还能直接写成 $w_t=1.5-0.5\times0.6^t,b_t=1-0.4^t$。

\textbf{因子绝对值小于 1，参数偏差就逐步缩小。}
解 $|r_w|<1$ 得 $0<\eta<1.5$，
解 $|r_b|<1$ 得 $0<\eta<1$；两者同时收敛的区间是 $\boxed{0<\eta<1}$。
负因子意味着在最低点两侧交替移动；绝对值仍小于 1 时，震荡会逐步缩小。
\begin{center}\small
\tagpdfsetup{table/header-rows={1}}
\begin{tabularx}{\linewidth}{@{}rrrX@{}}
\toprule
$\eta$ & $r_w$ & $r_b$ & 从本例初值出发的结果\\\midrule
0.3 & 0.6 & 0.4 & 两个方向单调靠近最优点\\
0.75 & 0 & $-0.5$ & $w$ 一步到 $1.5$，$b$ 震荡收敛\\
1 & $-1/3$ & $-1$ & $b$ 在 0 与 2 之间等幅摆动\\
1.2 & $-0.6$ & $-1.4$ & $w$ 收敛，$b$ 发散\\
1.5 & $-1$ & $-2$ & $w$ 等幅摆动，$b$ 发散\\
1.6 & $-17/15$ & $-2.2$ & 两个方向都发散\\
\bottomrule
\end{tabularx}
\end{center}
\textbf{都能收敛，也不是步子越大越快。}
仍从 $(1,0)$ 出发更新 10 步，$\eta=0.75$ 时损失约为 $0.500001$，
$\eta=0.9$ 时却约为 $0.511529$。两者都在稳定区间内，
但后一种的 $b$ 偏差每步乘 $-0.8$，比前一种的 $-0.5$ 缩小得慢。
因此要同时看步长、振荡和给定步数后的损失。

\begin{boundarybox}
这些阈值由本例损失的二次项系数决定，适用于恒定学习率和全量梯度。
换数据、损失尺度或优化器后，要重新判断步长；
某方向若从最优值开始，其参数偏差始终为零。
\end{boundarybox}

\begin{example}[breakable=false,title={{\symboliconfont ★} 选读证明：为什么这是全局最小？}]
\phantomsection\label{proof:three-point-global-minimum}
令梯度为零找到了候选点。将\eqref{eq:three-point-residual-loss} 展开、合并，再配方：
\begin{equation}
\begin{aligned}
L(w,b)&=\frac{2w^2-6w+3b^2-6b+9}{3}\\
&=\frac23\left(w-\frac32\right)^2+(b-1)^2+\frac12.
\end{aligned}
\label{eq:three-point-loss-surface}
\end{equation}
两个平方项都非负，在 $(w,b)=(3/2,1)$ 时同时为零。
这就证明了：全部参数中的最小损失恰为 $1/2$，前面的候选点确实是全局最优解。
\end{example}

\subsection{样本多时，为什么每次只用一部分数据？}
\label{intro:18c-training-engineering}

前面每一步都用全部三个样本求梯度，叫作\textbf{全批次梯度下降}。
样本很多时，每走一步都遍历全部数据，计算量可能很大。
可以每次只读一部分样本，用它们估计整体梯度，再更新一次参数；
这样减少了每次更新需要处理的数据量，也会让更新方向更受本次样本的影响。

\subsubsection{同一组数据，全量梯度与单样本梯度有何不同？}
\phantomsection\label{experiment:sample-gradient}

仍在 $(w,b)=(1,0)$ 处。先区分一条样本的损失与全量平均：
\[
\ell_i(w,b)=(wx_i+b-y_i)^2,\qquad L=\frac{\ell_1+\ell_2+\ell_3}{3}.
\]
例如样本 $(1,2)$ 有 $\ell=(w+b-2)^2$。一次只改 $w$ 或只改 $b$，
括号内的数值都会随它等幅变化。因此在当前参数处
\[
\frac{\partial\ell}{\partial w}=2(w+b-2)\cdot1=-2,\qquad
\frac{\partial\ell}{\partial b}=2(w+b-2)\cdot1=-2.
\]
单样本梯度不除以 $3$；将三条梯度取平均时才除以 $3$。
样本 $(0,2)$ 则有 $\ell=(b-2)^2$，梯度为 $(0,-4)$。
下表四行都在\textbf{同一组参数}处独立计算，用来比较不同的数据选择。
\begin{center}
\tagpdfsetup{table/header-rows={1}}
\begin{tabular}{lc}
\toprule
本次使用的数据 & 梯度（先 $w$，后 $b$）\\\midrule
只有样本 $(-1,-1)$ & $(0,0)$\\
只有样本 $(0,2)$ & $(0,-4)$\\
只有样本 $(1,2)$ & $(-2,-2)$\\
全部三个样本，梯度取平均 & $(-2/3,-2)$\\
\bottomrule
\end{tabular}
\end{center}
第一个样本已经预测正确，所以它给出的梯度为零；另外两个样本仍有误差。
使用全部样本，能综合它们的影响。
若等概率随机抽取一个样本，用对应行的梯度估计整体梯度并更新参数，
称为\textbf{单样本随机梯度下降}（Stochastic Gradient Descent，SGD）。
三个单样本梯度的平均恰好是全量梯度，但一次抽样只会得到其中一项。

\subsubsection{每次用几条样本？全部数据怎样分批训练？}

每次计算全部样本，计算量可能很大；上一小节又看到，只看一条样本时，
得到的梯度会受这条样本影响。能不能每次取几条，综合它们的影响，再更新一次？

仍在 $(w,b)=(1,0)$ 处，这次一起取原数据中的 $(0,2)$ 和 $(1,2)$ 两条样本。
由上表，它们的梯度分别为 $(0,-4)$ 和 $(-2,-2)$，平均后得到
\[
\frac{(0,-4)+(-2,-2)}{2}=(-1,-3).
\]
若仍取学习率 $\eta=0.3$，就用这个平均梯度更新一次：
\[
(w,b)=(1,0)-0.3(-1,-3)=(1.3,0.9).
\]
\phantomsection\label{concept:mini-batch}
这次一起使用的两条样本组成一个\textbf{小批次（Mini-batch）}。
\phantomsection\label{concept:batch-size}
每批包含多少条样本，叫作\textbf{批量大小（Batch Size）}，这里是 $2$。
把这一批样本的梯度取平均，再更新一次参数，称为\textbf{小批次 SGD}。
\textbf{同一批中的样本在同一组当前参数下求梯度；这一批更新后，下一批用新参数重新计算。}

样本更多时，怎样安排这些批次？为了看清分批流程，暂把数据规模扩展为 $100$ 条，
只讨论处理顺序，每批取 $10$ 条。
先随机打乱全部样本，取前 $10$ 条求平均梯度、更新一次；
再取接下来的 $10$ 条，用新参数求平均梯度、再更新一次。
这样依次处理，经过 $100/10=10$ 次参数更新，全部 $100$ 条样本就都处理了一遍。
\phantomsection\label{concept:epoch}
完整处理一遍训练数据，叫作一个\textbf{训练轮次（Epoch）}。
本例中，\textbf{每处理一批更新一次，处理完十批才完成一轮}，如图~\ref{fig:minibatch-epoch-cycle} 所示。
若继续训练下一轮，可以重新打乱这些样本，再依次分批处理。
\par\noindent\begin{minipage}{\linewidth}
\tagstructbegin{tag=Figure,alttext={一百条训练样本先随机打乱，再依次分成十批，每批十条。图中依次展示第一批、第二批和第十批，中间省略其余批次。每批用当前参数计算平均梯度并更新一次，箭头把更新后的新参数传给下一批。全部一百条样本处理一遍，完成一个训练轮次，累计更新十次。}}
\begin{center}
\begin{tikzpicture}[>=Latex,x=1mm,y=1mm,
 batch/.style={draw=linegray,rounded corners=2pt,fill=paper,
 align=center,font=\footnotesize,text width=35mm,inner sep=1.5mm,minimum height=18mm}]
\node[font=\small] at (55,16) {$100$ 条样本，先随机打乱，再依次分批};
\node[batch] (first) at (0,0) {第 $1$ 批：$10$ 条\\求平均梯度\\第 $1$ 次参数更新};
\node[batch] (second) at (50,0) {第 $2$ 批：$10$ 条\\求平均梯度\\第 $2$ 次参数更新};
\node[font=\small] (omitted) at (80,0) {$\cdots$};
\node[batch] (last) at (110,0) {第 $10$ 批：$10$ 条\\求平均梯度\\第 $10$ 次参数更新};
\draw[->,thick,flowcolor] (first.east)--node[above,font=\scriptsize] {新参数} (second.west);
\draw[->,thick,flowcolor] (second.east)--(omitted.west);
\draw[->,thick,flowcolor] (omitted.east)--(last.west);
\draw[thick,flowcolor] (-19,-13)--(-19,-16)--(129,-16)--(129,-13);
\node[font=\small] at (55,-23) {全部 $100$ 条处理一遍：$1$ 轮，累计更新 $10$ 次};
\end{tikzpicture}
\captionof{figure}{本例每批更新一次参数，全部数据处理一遍完成一个训练轮次}
\label{fig:minibatch-epoch-cycle}
\end{center}
\tagstructend
\end{minipage}\par


每批为什么取 $10$ 条？这是本例为说明流程所作的选择。
每批大一些，通常能减小抽样带来的梯度波动，但每次更新要处理更多样本；
每批小一些，单次更新处理的数据少，梯度估计的波动通常更明显。
小批次梯度是整体梯度的估计，更新方向会随本次样本变化，全量训练损失可能暂时上升。
这里的“随机”来自样本的抽取或打乱顺序。

回到原来的三点数据：每批取 $1$ 条时，一轮依次处理三条样本、更新 $3$ 次；
每批取 $3$ 条时，一次就用到全部样本，一轮更新 $1$ 次，正是前几节的全批次梯度下降。

\phantomsection\label{concept:optimizer}
按梯度修改参数的具体规则叫\textbf{优化器}，正文用上述梯度下降规则演示训练。
动量、Adam、梯度裁剪与训练不稳定的诊断要用到梯度向量和多步历史，
放在\hyperref[app:optimizer-details]{附录~\ref*{app:optimizer-details}（第~\pageref*{app:optimizer-details}~页）}逐步手算，学完向量后再读。\par\medskip

\chapterchecktasks{16-loss-gradient}{理解与复算}{
\item \textbf{预测：}沿用三点损失 $L(w,b)=\frac23(w-\frac32)^2+(b-1)^2+\frac12$，从 $(1,0)$ 出发。若学习率由 $0.3$ 改为 $0$，参数与损失会怎样变化？这为什么不等于“模型已收敛”？\chapterinlinehint{把零学习率代入更新式，再检查梯度。}
\item \textbf{手算：}写出 $(w,b)=(1,0)$ 的两个偏导数；以 $\eta=0.3$ 同步更新 $w,b$，再代回损失函数。\chapterinlinehint{先在旧参数处求梯度，再同步更新并复算损失。}
\item \textbf{找错：}先更新 $w$，再用新 $w$ 与旧 $b$ 求另一个梯度，为什么不算同一次同步更新？\chapterinlinehint{给本次迭代的所有参数标上同一个步数。}
}


\par\smallskip
本例可以直接求解，也适合核验梯度下降怎样从当前参数逐步接近最优解。
梯度下降提供可反复执行的更新过程；样本多时，小批次 SGD 用部分样本估计梯度，
减少每次更新的数据处理量，同时带来方向上的波动。
接下来研究：若函数只是一条直线，怎样表示弯曲的关系？
